% Original analytic teaching diagram, 2026-11-12.
% X = 1{U <= theta}, U ~ Unif(0,1), theta = 0.4, perturbed to theta + h = 0.55 (scale 90 mm per unit).
% Blue: {U <= theta}, where X = 1.  Red strip (theta, theta+h]: the only U whose X changes (0 -> 1); probability h.
% Green dots: the four teaching uniforms 0.37, 0.80, 0.15, 0.62; none lies in the strip, so each X is locally
% constant in theta and the pathwise derivative of every sample is 0, while J(theta) = theta has derivative 1.
\definecolor{VB}{HTML}{7188A3}\definecolor{VBF}{HTML}{E5EAF0}
\definecolor{VG}{HTML}{829C8B}\definecolor{VGF}{HTML}{E8EEE9}
\definecolor{VY}{HTML}{B4A16F}\definecolor{VYF}{HTML}{F1ECDD}
\definecolor{VR}{HTML}{A7787D}\definecolor{VRF}{HTML}{F0E5E6}
\definecolor{VP}{HTML}{9588A5}\definecolor{VPF}{HTML}{ECE8F0}
\definecolor{VI}{HTML}{3B4252}\definecolor{VM}{HTML}{6B7280}\definecolor{VL}{HTML}{4C4D4F}
\begin{tikzpicture}[x=90mm,y=1mm,>=Stealth,font=\figurefont\fontsize{8}{10}\selectfont,
lab/.style={text=VM,font=\figurefont\fontsize{7.5}{9}\selectfont}]
\fill[VBF] (0,0) rectangle (0.4,8);
\draw[VB,line width=.7pt] (0,0) rectangle (0.4,8);
\fill[VRF] (0.4,0) rectangle (0.55,8);
\draw[VR,line width=.7pt] (0.4,0) rectangle (0.55,8);
\draw[VL,line width=.5pt] (0.55,0) rectangle (1,8);
\node[text=VI] at (0.2,6) {$X=1$};
\node[text=VI] at (0.775,6) {$X=0$};
\node[lab,text=VR,anchor=south] at (0.475,8.6) {$X$ 由 $0$ 变 $1$};
\draw[VL,line width=.5pt] (0,0) -- (1,0);
\foreach \u/\t in {0/$0$,0.4/$\theta$,0.55/$\theta+h$,1/$1$} {
  \draw[VL,line width=.5pt] (\u,0) -- (\u,-1.2);
  \node[text=VI,anchor=north] at (\u,-1.6) {\t};
}
\foreach \u in {0.37,0.80,0.15,0.62} {
  \fill[VG] (\u,2.3) circle (1mm);
}
\end{tikzpicture}
